\chapter{ラグランジアン、場、局所粒子過程}

ボソン場、フェルミオン場、電磁場はいずれも波動型の方程式に従う。それらの方程式には、場そのものとその空間微分および時間微分が含まれる。各場ごとに別々の運動方程式を列挙することもできるが、ラグランジアンは系全体を一つの簡潔な式にまとめる。その式にオイラー--ラグランジュの規則を適用すると、すべての場の方程式が一度に得られる。

この構成は停留作用と密接に関係している。ここでの当面の課題はより具体的である。場のラグランジアンが何に依存するかを特定し、それを使って古典的な方程式をいくつか求め、次に同じ場の式を量子力学的に再解釈する。量子論では、ラグランジアンが粒子の運動と相互作用を支配し、最終的には衝突の振幅と断面積を支配する。

\section{場の方程式のための簡潔なコード}

一つのスカラー場について、ラグランジアン密度は場とその一階微分に依存する。
\begin{equation}
\mathcal L
=
\mathcal L\!\left(\phi,\partial_t\phi,\partial_x\phi,\ldots\right)
=
\mathcal L\!\left(\phi,\partial_\mu\phi\right).
\end{equation}
時間微分と空間微分を明示的に列挙した形は計算中にはしばしばより有用であり、表記\(\partial_\mu\phi\)は同じ列挙を相対論的な形にまとめたものである。

\begin{figure}[htbp]
\centering
\includegraphics[width=0.82\textwidth]{lecture_09_figure_03.png}
\caption{スカラー場のラグランジアンの場および時空微分への依存性。}
\lectureframe{00:03:56}
\end{figure}

場のオイラー--ラグランジュ方程式は
\begin{equation}
\frac{\partial}{\partial t}
\frac{\partial\mathcal L}{\partial\dot\phi}
+
\frac{\partial}{\partial x}
\frac{\partial\mathcal L}{\partial\phi_x}
+
\frac{\partial}{\partial y}
\frac{\partial\mathcal L}{\partial\phi_y}
+
\frac{\partial}{\partial z}
\frac{\partial\mathcal L}{\partial\phi_z}
=
\frac{\partial\mathcal L}{\partial\phi},
\end{equation}
ここで
\begin{equation}
\dot\phi=\frac{\partial\phi}{\partial t},
\qquad
\phi_x=\frac{\partial\phi}{\partial x},
\qquad
\phi_y=\frac{\partial\phi}{\partial y},
\qquad
\phi_z=\frac{\partial\phi}{\partial z}.
\end{equation}
である。\(\mathcal L\)が複数の場を含む場合、この規則は各場について一度ずつ適用される。その結果が運動方程式の系である。

最も単純なスカラーの例は
\begin{equation}
\begin{split}
\mathcal L
&=
\frac12\dot\phi^2
-\frac12\phi_x^2
-\frac12\phi_y^2
-\frac12\phi_z^2
-V(\phi)
\\
&=
\frac12\dot\phi^2-\frac12(\boldsymbol\nabla\phi)^2-V(\phi).
\end{split}
\end{equation}
である。空間微分の項は時間微分の項と反対の符号を持つ。項ごとに微分すると
\begin{align}
\frac{\partial\mathcal L}{\partial\dot\phi}&=\dot\phi,
&
\frac{\partial\mathcal L}{\partial\phi_i}&=-\phi_i,
&
\frac{\partial\mathcal L}{\partial\phi}&=-V'(\phi),
\end{align}
となり、したがって
\begin{equation}
\ddot\phi-\nabla^2\phi=-V'(\phi).
\end{equation}
となる。

\begin{figure}[htbp]
\centering
\includegraphics[width=0.74\textwidth]{lecture_09_figure_04.png}
\caption{オイラー--ラグランジュ方程式、スカラー場のラグランジアン、および\(V(\phi)\)が生成する場の方程式。}
\lectureframe{00:08:23}
\end{figure}

これがスカラーラグランジアンから生成される方程式である。\editorialnote{保持されたフレームは左端で切り取られている。表示されている補完は、すぐ周囲の口述による導出によって確定される。}

ここで\(V(\phi)\)は場のポテンシャルエネルギー密度である。\(V\)への定数の寄与は、その微分が消えるため、この方程式に影響を与えない。線形ポテンシャルは代数的には可能であるが、下に有界でない。場のエネルギーは負の無限大に向かって流れうる。したがって最初の安定な例は
\begin{equation}
V(\phi)=\frac12m^2\phi^2,
\qquad
V'(\phi)=m^2\phi,
\end{equation}
であり、これは
\begin{equation}
\ddot\phi-\nabla^2\phi=-m^2\phi.
\end{equation}
を与える。

\(m\)の意味を明らかにするために、平面波を代入する。
\begin{equation}
\phi(\mathbf x,t)=\phi_0e^{i\omega t-i\mathbf k\cdot\mathbf x}.
\end{equation}
すると
\begin{equation}
\ddot\phi=-\omega^2\phi,
\qquad
\nabla^2\phi=-\mathbf k^2\phi,
\end{equation}
となるので、波動方程式は
\begin{equation}
-\omega^2\phi+\mathbf k^2\phi=-m^2\phi.
\end{equation}
となる。\(\phi\)を消去すると、
\begin{equation}
\omega^2=\mathbf k^2+m^2.
\end{equation}
となる。\(\hbar=c=1\)とする単位では、振動数はエネルギーであり、波数は運動量である。したがって
\begin{equation}
E^2=\mathbf p^2+m^2,
\end{equation}
となり、\(c\)の冪を復元すると、\(E^2=\mathbf p^2c^2+m^2c^4\)となる。パラメータ\(m\)は場の一量子の質量である。

相対論的不変性は、簡潔な式のレベルで課すことができる。\(\mathcal L\)がローレンツスカラーであれば、その運動方程式はローレンツ不変である。時間微分と空間微分の間の相対的な符号は、特殊相対性理論でおなじみの時空の符号である。

\section{非線形性と結合した古典場}

場のより高次の冪は新しい種類の振る舞いをもたらす。もし
\begin{equation}
V(\phi)=\frac12m^2\phi^2+g\phi^3,
\end{equation}
ならば
\begin{equation}
\ddot\phi-\nabla^2\phi=-m^2\phi-3g\phi^2.
\end{equation}
である。ラグランジアンにおいて場について二次の項は場について線形な方程式を与える。三次以上の項は非線形方程式を与える。

この区別は波束の運動に見て取れる。線形方程式では、右向きの波束と左向きの波束を加えて別の解を作ることができる。それらはしばらく重なり干渉した後、互いを通り抜け、元の形のまま現れる。非線形項があると、波束は互いに影響を及ぼす。それらは散乱し、変形し、あるいはより複雑な波に分裂しうる。線形波の一時的な干渉は相互作用ではない。波束の持続的な変化が非線形性の印である。

同じパターンは二つのスカラー場でも現れる。
\begin{align}
\mathcal L
&=
\frac12(\partial_\mu\phi)(\partial^\mu\phi)
+\frac12(\partial_\mu\rho)(\partial^\mu\rho)
-V(\phi,\rho),
\\
V(\phi,\rho)
&=
\frac12m^2\phi^2+\frac12M^2\rho^2+\rho\phi^2.
\end{align}
としよう。最後の項がなければ、二つの方程式は独立であり、
\begin{equation}
\ddot\phi-\nabla^2\phi=-m^2\phi,
\qquad
\ddot\rho-\nabla^2\rho=-M^2\rho.
\end{equation}
である。それらは質量\(m\)と\(M\)を持つ二種の量子を記述する。混合項を含めると、
\begin{align}
\ddot\phi-\nabla^2\phi&=-m^2\phi-2\rho\phi,
\\
\ddot\rho-\nabla^2\rho&=-M^2\rho-\phi^2.
\end{align}
となる。各方程式は今や他方の場を含んでいる。方程式は結合しており非線形である。\(\phi\)の波束は\(\rho\)の波束を散乱し、またその逆も成り立つ。微分項はエネルギー--運動量関係を支配し、微分を含まない二次の項が質量を定め、混合した高次の項が相互作用を表す。これが基本的な古典的パターンである。

\section{ディラック場とスカラー結合}

ディラック場では、方程式は時間について1階である。1次元の略記では、3つの空間方向への拡張は了解済みとして、その項構造は
\begin{equation}
i\dot\Psi
\simeq
i\alpha\frac{\partial\Psi}{\partial x}+m\beta\Psi,
\qquad
i\boldsymbol\alpha\cdot\boldsymbol\nabla\Psi
=i\sum_i\alpha_i\frac{\partial\Psi}{\partial x_i}.
\end{equation}
である。行列\(\alpha_i\)と\(\beta\)は\(\Psi\)の成分に作用する。対応するラグランジアンには、\(\partial_t\Psi\)に線形な項、\(\alpha_i\)を含む空間微分の双線形項、そして質量双線形項\(m\Psi^\dagger\beta\Psi\)が含まれる。ここで必要なのは、この構造上の事実だけである。\editorialnote{試みられたディラックラグランジアンにおける符号、括弧、および\(i\)の配置は、原典において明示的に未確定のままにされている。そこから規約に依存する厳密な公式が推論されることはない。}

\begin{figure}[htbp]
\centering
\includegraphics[width=0.76\textwidth]{lecture_09_figure_05.png}
\caption{ディラック場の議論における時間微分項、空間微分項、質量項。}
\lectureframe{00:25:37}
\end{figure}

\begin{classroomqa}
\lecturetimestamp{00:26:37}
\audiencequestion{試みられたディラックラグランジアンにおいて、括弧は正しく配置されており、\(i\)が式全体に掛かっているのか。}
\lecturerresponse{ここではその配置を確信をもって確定しておらず、また構造上の要点、すなわちディラック方程式を生成するラグランジアンが存在するという点には必要ない。}
\end{classroomqa}

ここでスカラー場\(\phi\)を加える。素朴な積\(\Psi^\dagger\Psi\,\phi\)はローレンツスカラーではない。質量項に現れるのと同じディラック行列を挿入すると
\begin{equation}
\mathcal L_{\mathrm{int}}
\propto
\Psi^\dagger\beta\Psi\,\phi.
\end{equation}
が得られる。等価に、\(\Psi^\dagger\beta\Psi\)がスカラーディラック双線形である。完全なラグランジアンを変分すると、ディラック方程式に\(\phi\)に比例する項が加わり、スカラー方程式に\(\Psi^\dagger\beta\Psi\)に比例する源が加わる。2つの場は互いに散乱する。

\begin{classroomqa}
\lecturetimestamp{00:28:03}
\audiencequestion{スカラー--フェルミオン相互作用に挿入される\(\beta\)は、ディラック質量項に現れるのと同じ\(\beta\)か。}
\lecturerresponse{はい。同じディラック行列である。}
\end{classroomqa}

比較
\begin{equation}
m\Psi^\dagger\beta\Psi
\qquad\hbox{and}\qquad
\Psi^\dagger\beta\Psi\,\phi
\end{equation}
により、結合は質量パラメータがスカラー場に置き換えられたもののように見える。この観察は、非零の真空値をもつ場を終盤近くで考察する際に再び現れる。

これまでの解釈は古典的であった。すなわち、波束が他の波束を散乱する。同じ代数は、場が演算子になると、より鮮明な粒子としての意味をもつ。

\section{量子場の冪は何をするか}

実スカラー量子場には生成演算子と消滅演算子の両方が含まれる。模式的には、
\begin{equation}
\phi(x)\sim A^+(x)+A^-(x),
\end{equation}
ここで \(A^+\) は量子を生成し、\(A^-\) は量子を消滅させる。したがって、\(\phi^3(x)\) には四種類の局所過程が含まれる：
\begin{equation}
\phi^3
\sim
(A^+)^3+(A^+)^2A^-+A^+(A^-)^2+(A^-)^3,
\end{equation}
ここで順序の詳細と数値係数は、この模式的読みでは省略されている。四つの粒子数の変化は
\begin{equation}
3\longrightarrow0,
\qquad
0\longrightarrow3,
\qquad
2\longrightarrow1,
\qquad
1\longrightarrow2.
\end{equation}
三つの場の因子はすべて同じ点 \(x\) で作用する。したがって、一つの三次の項には、三粒子消滅、三粒子生成、二つの量子の一つへの融合、および一つの量子の二つへの分裂が含まれる。物理的振幅では、これらの局所的な読みの一部は後にエネルギー--運動量保存によって除去されるかもしれないが、それらはすべて演算子積自体の中に存在する。

二次の自由項は異なる役割をもつ。まず、隣接する点 \(x\) と \(x'\) の差の規格化を抑制する：
\begin{equation}
\partial_x\phi
\ \propto\
\phi(x)-\phi(x').
\end{equation}
二乗して展開すると、項構造
\begin{align}
\frac12(\partial_x\phi)^2
&\ \propto\
\frac12\bigl[\phi(x)-\phi(x')\bigr]^2
\\
&=
\frac12\phi(x)^2
+\frac12\phi(x')^2
-\phi(x)\phi(x').
\end{align}
が得られる。各同一点の二乗にはいくつかの演算子の読みがある：それは二つの量子を生成するか、二つを消滅させるか、またはその同じ点で一つを消滅させ再生成するかもしれない。これらの項は粒子を輸送しない。交差項は \(x'\) で量子を消滅させ、\(x\) でそれを生成するか、またはその逆を行うことができる。したがって、それは最近接ホッピング項として作用する。繰り返される因子は粒子を \(x\) から \(x'\) へ、次に \(x''\) へ、そしてさらに先へと運ぶ。二次の微分項は自由伝播を支配し、三次以上の項は粒子を散乱し、生成し、消滅させる。

\begin{classroomqa}
\lecturetimestamp{00:36:07}
\audiencequestion{微分の有限差分表示には間隔 \(x-x'\) を含めるべきか？}
\lecturerresponse{はい。間隔を \(L\) と呼ぶことで、差を二乗した後に因子 \(L^{-2}\) が与えられる。その数値因子は係数を変えるが、ホッピングの解釈は変えない。}
\end{classroomqa}

間隔を復元すると
\begin{align}
\frac12(\partial_x\phi)^2
&\longrightarrow
\frac{1}{2L^2}\bigl[\phi(x)-\phi(x')\bigr]^2
\\
&=
\frac{1}{2L^2}\phi(x)^2
+\frac{1}{2L^2}\phi(x')^2
-\frac{1}{L^2}\phi(x)\phi(x').
\end{align}
が得られる。

質量項も二次であるが、微分を含まない：
\begin{equation}
\frac12m^2\phi(x)^2.
\end{equation}
その一粒子の読みは、入射量子が同じ点で吸収され再放出されることである。それは量子を動かさない；そこでのその存在を重み付けし、あるいはその意味で数え上げる。同じ演算子の二乗には、その二生成および二消滅の読みも含まれる。したがって、運動に関連する区別は、隣接する点を結ぶ交差項と、一つの点内で作用する質量挿入との間にある。

\section{量子電磁力学の局所頂点}

これらの素粒子過程の略図は、ファインマン図の出発点である。中心となる例は量子電磁力学、すなわち電子、陽電子、光子の理論である。

ディラック方程式には正エネルギーと負エネルギーのモードがある。古い言葉では、\(\Psi\) は両方の種類のモードの消滅演算子を含み、\(\Psi^\dagger\) は対応する生成演算子を含む。負エネルギーの電子を取り除くことは、正エネルギーの陽電子を生成することとして再解釈される。そのように読み替えた後、有用な模式的分解は
\begin{align}
\Psi
&\sim
C^-_{e^-}+C^+_{e^+},
\\
\Psi^\dagger
&\sim
C^+_{e^-}+C^-_{e^+}.
\end{align}
となる。したがって \(\Psi\) は電子を消滅させるか陽電子を生成し、\(\Psi^\dagger\) は電子を生成するか陽電子を消滅させる。この略記には規格化、スピンラベル、運動量積分のいずれも意味されていない。

QED相互作用はディラック電流を電磁ベクトルポテンシャルに結合させる。ここで用いる成分表記では、その構造は
\begin{equation}
\mathcal L_{\mathrm{QED,int}}
\sim
e\left(
\Psi^\dagger\Psi\,A_0
+
\Psi^\dagger\boldsymbol{\alpha}\Psi\mathbin{\boldsymbol{\cdot}}\mathbf A
\right),
\end{equation}
であり、ここで
\begin{equation}
J^0=\Psi^\dagger\Psi,
\qquad
\mathbf J=\Psi^\dagger\boldsymbol{\alpha}\Psi.
\end{equation}
である。記号 \(\sim\) は重要である。時間項と空間項の間の相対符号は計量とベクトルポテンシャルの規約に依存する。不変な言明は、四元電流が四元ポテンシャルと縮約されるということである。\editorialnote{出典は電流の時間成分と空間成分を定めるが、成分の縮約に対する計量の規約は定めていない。}

光子場\(A_\mu\)自体には、光子の生成演算子と消滅演算子が含まれている。そのベクトルとしての性質はまた、光子の偏光---方向を示す小さな``旗''を、その運動量に加えて担っている。

\begin{classroomqa}
\lecturetimestamp{00:44:28}
\audiencequestion{\(A_\mu\) がベクトルであり、フェルミオン因子が最初はスカラーのように見えるのに、なぜ相互作用がスカラーになりうるのか。}
\lecturerresponse{フェルミオンもまた四元ベクトルを与える。\(\Psi^\dagger\Psi\) はその時間成分であり、\(\Psi^\dagger\boldsymbol\alpha\Psi\) はその空間成分である。その電流を \(A_\mu\) と縮約するとスカラーが生じる。}
\end{classroomqa}

\(\Psi^\dagger\)、\(\Psi\)、\(A_\mu\) の演算子内容を掛け合わせると、一つの局所頂点について四つのよく知られた読みが得られる。
\begin{enumerate}
\item 電子が光子を放出する。
\item 電子が光子を吸収する。
\item 電子と陽電子が対消滅して光子になる。
\item 光子が電子--陽電子対を生成する。
\end{enumerate}
第一の読みでは、\(\Psi\) が入射電子を消滅させ、\(\Psi^\dagger\) が出射電子を生成し、光子場が光子を生成する。光子の消滅の部分を選べば吸収が得られる。陽電子の部分を選べば対消滅または対生成が得られる。一つの局所項がこの族全体を含んでいる。

フェルミオン流の図では、陽電子線上の矢印は陽電子の伝播方向と逆向きを指す。この規約により、フェルミオンの矢印は許されるすべてのQED頂点を連続的に通過する。四つの読みは、フェルミオン流の線を切るのではなく、外部脚を入射側と出射側の間で移すことによって異なる。

これらの読みのいずれも同じ係数 \(e\)、すなわち用いられている微視的規約における電荷を伴う。この係数は電荷 \(e\) であり、微細構造定数ではない。ただし両者は密接に関係している。頂点振幅は \(e\) に比例し、すべての運動学的因子と規格化因子を含めた後、対応する確率の尺度は \(e^2\) に比例する。

\begin{classroomqa}
\lecturetimestamp{00:49:44}
\audiencequestion{光子から電子--陽電子への読みは、頂点の他の読みと同じ係数をもつか。}
\lecturerresponse{はい。それは同じ電荷の係数 \(e\) を伴う。}
\end{classroomqa}

\begin{classroomqa}
\lecturetimestamp{00:49:59}
\audiencequestion{確率には単位がないのだから、スケーリング因子が必要なのではないか。}
\lecturerresponse{用いられている自然な微視的規約では、電荷 \(e\) は無次元である。それが振幅の尺度であり、その二乗が確率の尺度を与える。}
\end{classroomqa}

具体的な描像は、電子が壁に当たり急に止まるというものである。エネルギーと運動量が光子の形で去っていく確率がいくらかある。その放射振幅における素朴な電磁因子は \(e\) であるから、その確率の尺度は \(e^2\) を含む。この言明は結合因子を切り出したものであり、実際の確率は電子の状態と運動学にも依存する。

\begin{figure}[htbp]
\centering
\includegraphics[width=0.78\textwidth]{lecture_09_figure_06.png}
\caption{\(\Psi\) と \(\Psi^\dagger\) の電子および陽電子の演算子内容を、電子--光子相互作用の略図とともに示す。}
\lectureframe{00:50:39}
\end{figure}

\begin{classroomqa}
\lecturetimestamp{00:51:15}
\audiencequestion{電荷には単位があるのか。}
\lecturerresponse{ここで用いている自然な微視的慣習では、単位はない。巨視的な電荷はしばしばクーロンで表されるが、素粒子理論では電荷を選ばれた基本単位に対する相対量として測る。プランク単位系その他の自然単位系では、量は純粋な数として表される。}
\end{classroomqa}

この対比は慣習と尺度の問題である。クーロンは導線中の電流に便利であり、無次元の結合定数は微視的振幅に便利である。\editorialnote{00からの短い質問：52：37--00：53：15、``電荷を担う項の''まとめ方}については、かなり損なわれており、明確な構成は得られなかった。

初等的なQED相互作用における三つの場はすべて一つの時空点で評価される。頂点自体には内部的な時間遅延は割り当てられない。

\begin{classroomqa}
\lecturetimestamp{00:53:19}
\audiencequestion{電子の吸収と再放出、および光子の放出は同時なのか、それとも方程式には遅延が組み込まれているのか。}
\lecturerresponse{初等的な頂点では、すべての演算子が同じ時空点で評価される。したがって基本過程には、それらの作用の間に遅延は含まれない。}
\end{classroomqa}

\begin{classroomqa}
\lecturetimestamp{00:53:57}
\audiencequestion{それぞれの保存量は別々に保存されなければならないのか。}
\lecturerresponse{はい。電荷は別々に保存されなければならない量の一つである。}
\end{classroomqa}

電荷の保存は図にも代数にも現れる。入射電子を出射陽電子に変える過程は電荷を反転させるので、許されるQED頂点ではない。時間を上向きに描くと、電子の矢印は時間に沿って上向きであり、陽電子の矢印は下向きである。有効な頂点では、常に貫通するフェルミオン矢印がある。代数的には、相互作用は一つの\(\Psi\)と一つの\(\Psi^\dagger\)を含む。前者は電子を消滅させるか陽電子を生成するときに電荷を一単位上げ、後者はそれを一単位下げる。それらの効果は相殺する。

\begin{classroomqa}
\lecturetimestamp{00:56:05}
\audiencequestion{枠で囲まれた電子--光子の式がQEDラグランジアンの全体なのか。}
\lecturerresponse{いいえ。それは相互作用項である。伝播を支配する二次の自由光子項と二次のディラック項は書かれていない。}
\end{classroomqa}

\begin{classroomqa}
\lecturetimestamp{00:56:35}
\audiencequestion{古典的ラグランジアンは演繹され、量子ラグランジアンは推測され検証されるのか。また、運動エネルギー引くポテンシャルの構造は量子力学でも保たれるのか。}
\lecturerresponse{運動エネルギー引くポテンシャルのパターンは保たれる。時間微分が運動エネルギーの形を与え、空間勾配がそれを一般化し、微分されない場の項がポテンシャルの形を与える。量子ラグランジアンはより基本的な記述である。非常に多くの量子を含む状態では、同じラグランジアンが古典的波動方程式を生み出す。候補となるラグランジアンは物理的原理とより深い理論を用いて提案され、その後、衝突実験によって確立されるか棄却される。}
\end{classroomqa}

大きな占有数の極限は、同じ式が量子粒子論とその古典的波動の両方を支配できる理由を説明する。この一致は後から課された別個の奇跡ではない。古典場は多数の量子の集団的振る舞いである。逆に、実験的な衝突は原始的頂点を直接的に露わにし、QEDと標準模型の相互作用を非常に高い精度で決定することを可能にする。

\section{交差と時空の保存則}

局所頂点の脚は、入射位置と出射位置の間で交差させることができる。出射粒子を入射側へ移すと、それは反粒子として再解釈され、逆もまた同様である。この意味で、同じ電子--陽電子--光子演算子は、放出、吸収、対消滅、対生成、さらにはすべて入射という局所的な読み方さえもつ。

\begin{classroomqa}
\lecturetimestamp{00:59:34}
\audiencequestion{基本頂点図を回転させても、それは有効な局所的読み方であり続けるのか。}
\lecturerresponse{はい。その外部脚は入射位置と出射位置の間で移すことができ、それに対応する粒子--反粒子の再解釈を伴う。}
\end{classroomqa}

局所演算子の読み方は、まだ完全な物理過程ではない。その振幅は、頂点が起こり得るあらゆる時空点にわたって積分されなければならない。時間積分はエネルギー保存を課し、空間積分は運動量保存を課す。電子、陽電子、光子がすべて真空へ消えるという描像は、形式的な局所的読み方としては存在するが、その過程はエネルギーと運動量を保存できないため、積分された振幅は消える。

\begin{classroomqa}
\lecturetimestamp{01:00:30}
\audiencequestion{すべてが入射するという描像は、電子と陽電子と光子が消えて無になれるということを意味するのでしょうか。}
\lecturerresponse{それは局所演算子の読み方ですが、物理的な振幅には頂点位置に関する積分が含まれます。それらの積分がエネルギーと運動量の保存を課すため、この運動学的に禁じられた過程の振幅はゼロです。}
\end{classroomqa}

\begin{classroomqa}
\lecturetimestamp{01:01:57}
\audiencequestion{逆に真空から電子--陽電子--光子への読み方も存在するのでしょうか、また排除された過程は積分の後で消えるのでしょうか。}
\lecturerresponse{逆の局所的な読み方も存在します。空間と時間にわたる積分が、保存則と両立しないあらゆる過程を取り除きます。}
\end{classroomqa}

ラグランジアン中の微分は、生成消滅記号そのものを微分することなく演算子展開に適合します。場は模式的に次の形をとります。
\begin{equation}
\Psi(x)
\sim
\int d^3k\,
\bigl[
C(\mathbf k)e^{i\mathbf k\cdot\mathbf x-i\omega t}
+\cdots
\bigr].
\end{equation}
演算子\(C(\mathbf k)\)はこの時空展開における係数です。微分は平面波に作用し、\(\mathbf k\)または\(\omega\)を引き下ろします。

\begin{classroomqa}
\lecturetimestamp{01:02:27}
\audiencequestion{場の展開を微分するとき、生成演算子そのものが微分されるのでしょうか。}
\lecturerresponse{いいえ。生成消滅演算子はこの展開において時空座標に依存しません。微分は指数因子に作用し、運動量と振動数を引き下ろします。}
\end{classroomqa}

\section{有効理論としての核子とパイ中間子}

同じ言語で陽子、中性子、パイ中間子を記述できます。陽子と中性子はフェルミオンであり、場\(\Psi_p\)と\(\Psi_n\)で表されます。この相互作用の数え上げの議論では、パイ中間子の場は3つの電荷の種類をもつスカラー場として扱われます。
\begin{equation}
\pi^+,
\qquad
\pi^-,
\qquad
\pi^0.
\end{equation}
陽子と\(\pi^+\)は正の電荷1単位を担い、\(\pi^-\)は負の電荷1単位を担い、中性子と\(\pi^0\)は中性です。より深いレベルではパイ中間子はクォーク--反クォーク系ですが、その内部記述は最初は未解決のままにしておけます。

電荷を保存する核子--パイ中間子場の積の模式的なリストは次のとおりです。
\begin{equation}
\mathcal L_{N\pi}
\sim
\left\{
\Psi_p^\dagger\Psi_n\,\pi^+,\;
\Psi_n^\dagger\Psi_p\,\pi^-,\;
\Psi_p^\dagger\Psi_p\,\pi^0,\;
\Psi_n^\dagger\Psi_n\,\pi^0
\right\}.
\end{equation}
ここで示されているのは場の内容と電荷の割り当てだけであり、相対係数、符号、スピノル行列の構造はここでは固定されません。最初のパイ中間子場の電荷は、具体的な読み方を試すことで決まります。中性子と正のパイ中間子は陽子を生成しうるのに対し、そこに負のパイ中間子を用いると負の入射電荷が正の出射電荷に変わってしまいます。共役項は反対の電荷の過程を与えます。\editorialnote{パイ中間子の符号は、ここに示す電荷保存の対に到達するまでに資料中で何度か修正されており、示されている順序はその最終的な模式構造を記録しています。}

各項は再びいくつかの交差した読み方をもちます。出射中性子を入射側に移すと、それは入射反中性子に変わります。したがって荷電頂点は、陽子と反中性子が結合して\(\pi^+\)になることや、運動学が許すときには\(p\to n+\pi^+\)のような過程を記述できます。中性項は陽子または中性子が\(\pi^0\)を放出または吸収することを可能にします。電荷保存が許される族を選び、エネルギー、運動量、およびさらなる対称性があればそれらが、特定の実験でどの成員が起こりうるかを決定します。

これはクォークの下部構造が知られる前には有用な有効記述でした。探査がその構造を分解できないときにはいつでも有用であり続けます。陽子は構成要素の有限な分布をもち、ちょうど原子が内部構造をもつのに似ていますが、はるかに小さい尺度です。実験が陽子を分解したり分解能で見たりするには小さすぎる運動量しか移さないなら、陽子はそれ自身の場をもつ点粒子として扱えます。十分に硬い衝突はクォークを露わにし、そのとき点粒子としての陽子場はもはや適切ではありません。

陽子は荷電しているため、局所的な陽子--光子相互作用ももち、模式的には\(J_p^\mu A_\mu\)であり、ここで\(J_p^\mu\)は陽子の電流です。電子は対応する電子--光子相互作用をもちます。陽子が光子を放出して中性子に変わることはできません。\(p\to n+\gamma\)は電荷保存に反します。

\begin{classroomqa}
\lecturetimestamp{01:10:28}
\audiencequestion{陽子と電子が相互作用すると何が起こるのでしょうか。}
\lecturerresponse{各荷電粒子は光子頂点をもちます。陽子は陽子のままで光子を放出でき、電子は電子のままで光子を吸収できます。陽子--電子散乱は、それら2つの局所頂点と、その間の光子の伝播から組み立てられます。}
\end{classroomqa}

\section{局所性と繰り返される素過程}

陽子--電子の交換全体に対して、特別な非局所項は存在しない。一つの局所頂点が陽子において光子を放出し、光子の二次項がそれを隣接点から隣接点へと運び、別の局所頂点が電子においてそれを吸収する。この拡張された過程は、それら素過程の積である。
\begin{equation}
\begin{aligned}
\bigl[J_p^\mu A_\mu\bigr](x_p)
&\longrightarrow
\underbrace{x_p\to x_1\to x_2\to\cdots\to x_e}_{\text{quadratic photon propagation}}
\\
&\longrightarrow
\bigl[J_e^\nu A_\nu\bigr](x_e).
\end{aligned}
\end{equation}
この模式的な連鎖における矢印は、隣接点への挿入が逐次続くことを示すものであり、新たな非局所相互作用を意味するものではない。

\begin{classroomqa}
\lecturetimestamp{01:12:06}
\audiencequestion{陽子--電子の交換全体に対するラグランジアンは、どのように書けばよいのだろうか。}
\lecturerresponse{拡張された事象全体に対する項を付け加えるのではない。それは、局所的な陽子--光子頂点、繰り返される光子の二次伝播項、および局所的な電子--光子頂点から組み立てられる。}
\end{classroomqa}

局所構造には二つの形がある。未微分の相互作用項は一つの時空点で作用する。微分項は有限差分として表され、その点を隣接点とのみ結び付ける。どちらの種類も、粒子を有限距離にわたって直接移すことはない。長距離伝播は多数の無限小移動から組み立てられ、素過程としての遠隔作用は存在しない。

\begin{classroomqa}
\lecturetimestamp{01:14:36}
\audiencequestion{ラグランジアンは一点のみに項を含むのだろうか。}
\lecturerresponse{相互作用項は一点で作用する。微分項は一点を無限小に隣接する点と結び付ける。放出、多数の隣接ホップ、吸収が合わさって有限の交換を生み出す。}
\end{classroomqa}

\begin{classroomqa}
\lecturetimestamp{01:15:36}
\audiencequestion{拡張された過程の要素はすべて存在し、それらを合成することは掛け算を意味するのだろうか。}
\lecturerresponse{はい。個々の素過程である放出、ホップ、吸収はそれぞれラグランジアンに符号化されている。繰り返される挿入の積が、それらの要素をより大きな過程へと合成する。}
\end{classroomqa}

ここで議論している相対論的量子場の理論は、その場が定義される時空設定を前提としている。これは、動的な時空幾何学の完全な量子論と混同されるべきではない。

\begin{classroomqa}
\lecturetimestamp{01:16:28}
\audiencequestion{この合成は、ブラックホールの事象の地平面近傍のような強く相対論的な設定でも有効であり続けるのだろうか。}
\lecturerresponse{ブラックホールへの拡張は、この構成の範囲外である。ここで用いている量子場の理論は、それ自体では、重力が本質的となるブラックホールの動力学を完全に記述するものではない。}
\end{classroomqa}

場とは、素過程としての局所的な作用を表す演算子である。一つの場の因子は粒子を吸収することも、一つ放出することもできる。二つの因子の積は、同一の点で一つの粒子を吸収し別の粒子を放出することができ、より高次の積は粒子の内容をさらに変化させうる。局所演算子の積は、これらの作用を履歴へと結び付ける。時空において、運動は、ある小さな領域での吸収、隣接領域への再放出、さらなる吸収、さらなる再放出によって表される。ステップ数を増やしながらステップ幅をゼロに近づけると、連続的な軌道が得られる。

\begin{classroomqa}
\lecturetimestamp{01:17:30}
\audiencequestion{この局所ステップ描像において、場、距離、時間はどのように理解されるべきだろうか。}
\lecturerresponse{場は局所的な吸収と放出を表す演算子である。積はそれらの素過程を結び付け、時空を通る運動は、連続極限を取る前の無限小の隣接ステップから組み立てられる。}
\end{classroomqa}

したがって、正則化された出発点とは、時空を小さなセルに分割し、場の演算子をそれらのセルに付随させることである。\editorialnote{01における聴衆の発言は17:30--01:20:40の部分が一部不明瞭であり、言い換えた質問は復元可能な主題のみを保持しているが、局所演算子と時空の逐次ステップに関する応答は明確である。} 微分は差分になる。二次の交差項は隣接セルを結び付け、質量項と相互作用項は一つのセル内で作用する。セルは極限理論を定義するための道具であり、追加の粒子や新たな物理的質量スケールではない。

特殊相対論的効果を各ステップで別個に付け加える必要はない。ラグランジアンがローレンツスカラーであれば、そこから組み立てられる振幅はローレンツ不変であり、時間の遅れのような効果はその構造から従う。

\begin{classroomqa}
\lecturetimestamp{01:21:10}
\audiencequestion{高エネルギー運動と時間の遅れは、多段階過程において余分な問題を引き起こすのだろうか。}
\lecturerresponse{別個の補正は必要ない。ラグランジアンがローレンツスカラーであれば、時間の遅れのような特殊相対論的効果はすでに結果に組み込まれている。質問の背後にある正確な懸念は、それ以外の点では十分に明確ではない。}
\end{classroomqa}

生成演算子と消滅演算子は、位置基底で表すことも運動量基底で表すこともできる。位置空間の生成演算子は運動量空間の生成演算子のフーリエ重ね合わせであり、逆もまた同様である。位置と運動量の両方を正確に指定した状態を生成する演算子は存在しない。粒子を鋭く局在させるには広い帯域の運動量が必要であり、運動量で鋭く局在させると位置に広がる。二つの基底の間のフーリエ関係が、この言葉で表した不確定性原理である。

\begin{classroomqa}
\lecturetimestamp{01:22:48}
\audiencequestion{生成演算子と消滅演算子は、位置、運動量、および不確定性原理とどのように関係しているのか。}
\lecturerresponse{それらは位置の関数として、あるいは運動量の関数として表すことができる。二つの表現は互いのフーリエ変換である。鋭く局在した位置状態には多くの運動量が必要であり、鋭く局在した運動量状態は位置に広がる。\(x\) と \(p\) を同時に鋭くする生成演算子は存在しない。}
\end{classroomqa}

セルの構成は、これで完全に述べることができる。セルを \(n\) でラベル付けし、微分を例えば次のような有限差分に置き換え、
\begin{equation}
\partial_\mu\phi(x)
\longrightarrow
\frac{\phi_{n+\hat\mu}-\phi_n}{a},
\end{equation}
すべての場の種をすべてのセルに配置する。二次のクロス項は、一つのセルから量子を一つ取り除き、隣のセルにそれを生成する。連続理論は、正則化された表式が定義された後で間隔 \(a\) をゼロに近づけることによって得られる。
\begin{equation}
\underbrace{\phi_n\phi_{n+\hat\mu}}_{\text{hop }n\leftrightarrow n+\hat\mu},
\qquad
\underbrace{m^2\phi_n^2}_{\text{same-cell mass insertion}}.
\end{equation}

\begin{classroomqa}
\lecturetimestamp{01:25:16}
\audiencequestion{時空をセルに細分することを完成できるのか。}
\lecturerresponse{はい。場はセルによって添字付けられ、微分は有限差分になり、二次の項は粒子を隣接するセルの間でホップさせる。連続体は、セルの大きさがゼロに近づくにつれて回復される。}
\end{classroomqa}

電子はその二次の項を通じてセル間をホップし、光子も同様にする。電子と光子が一つのセルで出会うと、QED相互作用は光子を生成または消滅させながら電子を吸収して再放出することがある。スカラー質量項には同一セルでの読み方がある。粒子はそのセルに入り、\(m^2\) に比例する係数でそのセルを離れる。それはホップを与えない。

\begin{classroomqa}
\lecturetimestamp{01:28:18}
\audiencequestion{セルの大きさは質量係数と関係しているのか。}
\lecturerresponse{いいえ。セルの大きさは最終的にゼロにされる正則化の間隔であり、質量係数はラグランジアンにおける独立なパラメータである。}
\end{classroomqa}

\begin{classroomqa}
\lecturetimestamp{01:28:43}
\audiencequestion{局所因子を掛け合わせることは、単に時間的な順序ではなく、時空にわたる結合振幅の構成なのか。}
\lecturerresponse{はい。セルは時空を覆い、それらの局所的な寄与はその時空格子にわたって結合する。その積は、必要な局所事象と伝播の段階に対する結合振幅を与える。}
\end{classroomqa}

\section{源、検出器、およびそれらの間の力学}

完全な過程は時空における歴史である。その初期条件は実験者が準備するものによって定まり、最終条件は検出器が記録するものによって定まる。単純な源は電子を放出する熱フィラメントである。フィラメントにおける場の演算子がその入射電子を生成し、検出器の位置における場の演算子がその最終的な吸収を記述する。

\begin{classroomqa}
\lecturetimestamp{01:29:30}
\audiencequestion{そのような時空過程の初期段階とは何か。}
\lecturerresponse{それは源において準備される。熱フィラメントから放出された電子が単純な例である。フィラメントにおける場の演算子が電子を生成し、検出器における場の演算子が最終的な吸収事象を指定する。}
\end{classroomqa}

同じ構成は真空から述べることもできる。源における生成演算子が真空の上に入射状態を構成する。検出器に対応する消滅演算子が出射状態を選び出す。これらの演算子は境界条件を指定するのであり、源と検出器の間の領域全体にわたる力学を置き換えるものではない。

\begin{classroomqa}
\lecturetimestamp{01:31:11}
\audiencequestion{力学を系全体ではなく放出器のみに置くのは誤りだろうか。}
\lecturerresponse{はい。真空から始め、源の演算子を用いて入射粒子を準備し、検出器の演算子を用いて最終粒子を選び出す。それらの間では、繰り返されるラグランジアンの挿入が時空全体にわたって作用し、粒子を動かし、生成し、消滅させる。}
\end{classroomqa}

十分な数の挿入があれば、複数の粒子が多くの隣接するセルを通って同時に動くことができる。粒子は途中で生成されたり消滅したりすることもできる。初期の演算子は何が入るかを述べ、最終の演算子は何が記録されるかを述べ、繰り返されるラグランジアンの操作がそれらの間のすべてを与える。

\begin{classroomqa}
\lecturetimestamp{01:33:24}
\audiencequestion{任意個のラグランジアン挿入が起こりうるならば、それらの位置は積分されているのか。}
\lecturerresponse{はい。それらの時空位置は積分されている。ここでは厳密なファインマン則を導出しているのではない。当面の要点は、単純な局所場の表式が、完全な過程が組み立てられる元になる基本的相互作用を符号化しているということである。}
\end{classroomqa}

\section{位相対称性、電荷、振幅}

原始的な過程をラグランジアンに書き表すと、それらの対称性も可視化される。電子場の大域的位相変化を考えると、
\begin{equation}
\Psi\longrightarrow e^{i\theta}\Psi,
\qquad
\Psi^\dagger\longrightarrow e^{-i\theta}\Psi^\dagger,
\qquad
A_\mu\longrightarrow A_\mu.
\end{equation}
\(\Psi\)と\(\Psi^\dagger\)の数が等しい項は、位相が相殺するため不変である。数が等しくない項は位相を獲得するであろう。それらの数の等しさが、電荷保存の代数的な形である。

これはラグランジアン記述の中心的な利点である。その場の因子は素過程を表示し、それらの組み合わせは対称性を表示し、対称性は保存則を表示する。繰り返しの挿入は、用意された初期状態を選択された最終状態へと発展させる。その結果が量子振幅である。厳密な評価にはファインマン則が必要であるが、構造的な内容はすでに見えている。すなわち、局所的な操作は保存則によって許されるあらゆる次数で繰り返される。

\section{場の測定と交換子}

場の演算子は、演算子という通常の量子力学的な意味を保っている。その固有値は場の測定の可能な結果であり、その固有ベクトルはそれらの値に対応する状態を記述する。電場は、荷電粒子をその中に置き、その粒子がどのように加速されるかを観測することによって調べることができる。

\begin{classroomqa}
\lecturetimestamp{01:37:06}
\audiencequestion{場の演算子の固有値と固有ベクトルは、通常の量子力学的な役割を果たすのか。}
\lecturerresponse{はい。固有値は測定されうる場の値であり、場は粒子に対するその効果を通じて調べることができる。交換子は同時測定に対する制限を決定する。}
\end{classroomqa}

自明な恒等式
\begin{equation}
[\phi(x),\phi(x)]=0
\end{equation}
は測定の問題を解決しない。二つの観測量とそれらの時空関係を指定しなければならない。場とその共役運動量は正準的な非可換対をなし、時間的に関連する点における場の演算子は可換である必要はない。空間的に離れた点における局所ボソン場は、局所性が要求するように可換である。したがって、交換子と不確定性の間のよく知られた関係は引き継がれるが、それは演算子の対が正しく同定された後にのみである。\editorialnote{原著者は異なる点における非可換性を広く述べすぎている。ここでの限定は、その測定点を保ちつつ、空間的に離れた局所ボソン場の微視的因果性を課すものである。}

\begin{classroomqa}
\lecturetimestamp{01:39:57}
\audiencequestion{電場または電磁場はボルトで測定されるのか、それともボルト毎メートルで測定されるのか。}
\lecturerresponse{電場はボルト毎メートルで測定される。ボルトとメートルは微視的な場の過程にとって最も有用な言語ではなく、それらは光子の観点からより自然に整理される。逆に、量子場の理論は家に水をやる方法を決めるのに有用な言語ではない。}
\end{classroomqa}

ある記述に自然な単位が、別の記述において啓発的な単位であるとは限らない。ボルトとボルト毎メートルは巨視的な装置に適合し、光子の生成と消滅は微視的な過程に適合する。

\begin{classroomqa}
\lecturetimestamp{01:41:10}
\audiencequestion{中心的な考えは、ラグランジアンの項がより大きな過程へと組み合わされる構成要素であるということか。}
\lecturerresponse{はい。素過程はある点から次の点へと繰り返される。各基本要素は一つの時空点で作用するか、あるいは微分項を通じて隣接する点を結ぶ。見かけ上非局所的な伝播はこれらの局所的な段階から組み立てられており、遠隔作用の要素は存在しない。}
\end{classroomqa}

\section{局所過程の積と経路にわたる和}

局所過程の繰り返しを表示する簡潔な発見的方法がある。時空をセルに分割し、
\begin{equation}
\prod_x\bigl(1+\mathcal L_I(x)\bigr),
\end{equation}
と書く。ここで\(x\)はセルをラベルし、\(\mathcal L_I(x)\)は局所相互作用ラグランジアンである。これは厳密な連続体の公式ではない。これは、積がどのようにして局所事象のあらゆる許された選択を生成するかを明らかにする組み合わせ的な代用物である。

\begin{classroomqa}
\lecturetimestamp{01:42:18}
\audiencequestion{多数の生成、伝播、消滅の段階を一つの表式に圧縮できるか。}
\lecturerresponse{発見的には、各時空セルに対して一つの因子\(1+\mathcal L_I(x)\)を掛ける。積を展開すると、事象が無い場合、任意の選択された点における一つの事象、または複数の点における事象が生成される。}
\end{classroomqa}

2つのセルの場合、
\begin{align}
\bigl(1+\mathcal L_I(x)\bigr)
\bigl(1+\mathcal L_I(x')\bigr)
&=
1+\mathcal L_I(x)+\mathcal L_I(x')
\notag\\
&\quad{}+\mathcal L_I(x)\mathcal L_I(x').
\end{align}
最初の項は、どちらの相互作用も起こらないことを意味する。次の2つの項は、1つの素過程を\(x\)または\(x'\)に置く。最後の項は、両方の点に相互作用を置く。因子が多くなると、1点、2点、3点、…を含む項が現れる。いくつかの頂点をもつ項は、たとえば2つの入射粒子が3つの出射粒子を生成する過程を記述できる。

二次の伝播部分も同じ構成に入る。有限の変位は、隣接する挿入の連鎖によって表される。展開された積のどこかには、1つのセル、その隣、さらにその隣、…という因子があり、粒子を端点の間で運ぶ。

\begin{classroomqa}
\lecturetimestamp{01:46:29}
\audiencequestion{有限の伝播に多くの局所運動項が必要なら、必要な運動はどのようにして回復されるのか。}
\lecturerresponse{大きな積には、粒子を初期点から最終点まで運ぶ隣接挿入の連鎖が含まれる。一般に、異なる経路に沿ったそのような連鎖が多数ある。各経路は振幅に寄与し、全振幅はそれらの和である。}
\end{classroomqa}

\(r\)が経路をラベルすると、模式的には
\begin{equation}
\mathcal A_{\mathrm{total}}=\sum_r\mathcal A_r,
\qquad
P=\left|\mathcal A_{\mathrm{total}}\right|^2.
\end{equation}
通常の確率を経路に割り当ててから足し合わせるのではない。まず振幅を足し合わせる：
\begin{equation}
\left|\sum_r\mathcal A_r\right|^2
=
\sum_r|\mathcal A_r|^2
+
\sum_{r\ne s}\mathcal A_r\mathcal A_s^*.
\end{equation}
交差項が干渉である。

\begin{classroomqa}
\lecturetimestamp{01:48:27}
\audiencequestion{これは二重スリット実験と整合するか。}
\lecturerresponse{まさにそのとおり。2つのスリットを通る選択肢が振幅に寄与する。その和を2乗すると干渉項が生じる。ちょうど、局所ステップ構成における多くの時空経路が干渉するのと同じである。}
\end{classroomqa}

\begin{classroomqa}
\lecturetimestamp{01:49:00}
\audiencequestion{各因子\(1+\mathcal L_I\)に現れる\(1\)はどこから来るのか。}
\lecturerresponse{それは、セルが相互作用をしない寄与と相互作用をする寄与の両方に寄与できるようにする組合せ的な仕掛けである。経路積分の方法と関係するが、表示された積は技術的な表式へのヒューリスティックな近似にすぎない。}
\end{classroomqa}

厳密な形式化は、適切な時空積分、順序、規格化、位相を与える。この積が有用なのは、完全な形式化のふりをすることなく、組合せを明らかにするからである。

\begin{classroomqa}
\lecturetimestamp{01:49:57}
\audiencequestion{これは停留作用の原理か。}
\lecturerresponse{いいえ。古典的な停留作用は、特別な古典軌道を選び出す。量子力学は代わりに、作用によって重みが支配される軌道にわたって振幅を足し合わせる。非常に遠回りな経路---事務所から近くの家へ行く途中でパリを通るような経路---でさえ、量子力学的な和に含まれる。}
\end{classroomqa}

関係が微妙であっても、対比は単純である。古典力学は停留経路を選ぶ。量子力学は経路を足し合わせる。適切な極限では、位相が古典軌道の周りで強め合う。

\section{ラグランジアンの推論とセクターの選択}

ラグランジアンと指定された過程が与えられれば、その確率を計算する規則は確定している。計算には、難しい数学がしばしばそうであるように、工夫が必要かもしれないが、より深い不確実性は、自然がどの場と結合を用いるかを発見することにある。

\begin{classroomqa}
\lecturetimestamp{01:50:43}
\audiencequestion{芸術性はラグランジアンを表現することにあるのか、それとも使うことにあるのか。}
\lecturerresponse{主な芸術性は、実験データからラグランジアンを推論することにある。いったんそれが分かれば、2つの粒子が入り4つが出るような過程には、精密な規則によって計算される確定した振幅がある。数学を実行するには巧妙さが必要かもしれないが、それは理論の曖昧さではない。}
\end{classroomqa}

粒子衝突は素過程の頂点を露わにする。項が提案され、しばしば対称性とより深い原理に導かれて、得られる振幅が観測結果と照合される。成功すれば、許される場の組合せとその数値係数の両方が確定する。

\begin{classroomqa}
\lecturetimestamp{01:52:24}
\audiencequestion{相互作用ごとに別々のラグランジアンがあるのか---ラグランジアンの大きな本のようなものか。}
\lecturerresponse{作業上の描像は代わりに、すべての場と相互作用を含む1つの大きなラグランジアンである。ただし、その描像は最終的には誤りかもしれない。認識可能なセクターは和の部分として現れる。新しい粒子は場を追加し、新しい素過程は相互作用項を追加する。}
\end{classroomqa}

完全な表式は、本一冊の長さにならずとも、整然としないものになりうる。意図的に粗い目安として、素粒子と呼ばれる対象が百個程度、項が数百個程度を想像してもよい。量子電磁力学、強い相互作用と弱い相互作用、そして重力の部門を合わせても、ぎっしり書いた数ページしか占めないかもしれない。正確な数が要点ではない。作業上の描像は、多くの認識可能な部品から組み立てられた一つの動的な表式であり、それは明示的に、最終的には不完全であると判明するかもしれない描像である。

それらの場がより基礎的な時空幾何学を定めるかどうかは、ここで用いている標準的な量子場の枠組みの範囲を超えている。\editorialnote{01:55:00--01:55:20という問いは、多様体や幾何学的構成を信頼できる形で同定するには損傷が大きすぎる。確実なのは、この適用範囲の境界に関する応答だけである。}

\begin{classroomqa}
\lecturetimestamp{01:55:22}
\audiencequestion{実験的に決定された数は結合定数なのか。}
\lecturerresponse{結合定数はラグランジアン中の項の係数である。二次の質量項の係数\(m^2\)が一例であり、\(\phi^4\)、\(\phi^6\)、あるいはさらに高次の項の係数も他の例である。}
\end{classroomqa}

一つの係数が、関連する振幅の族を支配する。各場の積には複数の生成と消滅の読み方があるからである。したがって、ラグランジアンを知ることは、一つの孤立した反応よりもはるかに多くを決定する。

\begin{classroomqa}
\lecturetimestamp{01:56:05}
\audiencequestion{量子電磁力学を計算するには、ラグランジアン全体を知らなければならないのか。}
\lecturerresponse{無限の精度を求めるなら、電子と光子が相互作用しうるあらゆる部門が必要になる。しかし原子のエネルギーでは、通常、電子--光子の部門を切り離すことができる。微小なクォーク構造を分解するだけのエネルギーがないため、必要な精度の範囲で他の部門は無視できる。}
\end{classroomqa}

実際の計算では、実験のスケールに関係する部門を用いる。これは流体力学に似ている。複雑な微視的系が、そのより深い構成要素をすべて露わにすることなく、精密な粗視化記述を持ちうるのである。したがって量子場理論は、現象論的でありかつ厳密でありうる---実験から推測された場と係数に関しては現象論的であり、それらの材料が指定された後に適用される規則に関しては厳密である。

\begin{classroomqa}
\lecturetimestamp{01:58:09}
\audiencequestion{標準模型は重力を含むのか、またこれらの相対論的ラグランジアンは一般相対性理論を含むのか。}
\lecturerresponse{公式に標準模型と呼ばれる理論は重力を含まないが、より広い理論を論じる際には重力を付け加えてもよい。ここで考える場の理論は、量子力学と特殊相対性理論を組み合わせたものである。それ自体は一般相対性理論を含まない。}
\end{classroomqa}

\section{質量係数としての真空値}

スカラー場の質量は\(\phi^2\)の係数として現れる。したがって、別の場との相互作用は、その別の場が非零の真空値を持つとき、質量項を模倣しうる。

\begin{classroomqa}
\lecturetimestamp{01:59:16}
\audiencequestion{粒子の質量がヒッグス相互作用を通じて生じるなら、ヒッグス結合は質量項として働きうるのか。}
\lecturerresponse{はい。項\(H\phi^2\)は三次であり、通常は二本の\(\phi\)の脚と一本の\(H\)の脚を持つ相互作用を表す。ポテンシャル\(V(H)\)が非零の\(H\)を持つ真空を選ぶなら、その真空値は\(\phi^2\)に掛かる数であり、したがって有効な二次の質量項に寄与する。}
\end{classroomqa}

模式的に考えると
\begin{equation}
\mathcal L\supset H\phi^2-V(H).
\end{equation}
演算子の項として、\(H\phi^2\)は\(\phi\)を吸収し、\(\phi\)を放出し、\(H\)量子を放出または吸収しうる。ポテンシャルが\(H=0\)から離れた二つの極小を持ち、真空がその一方を選ぶとしよう。
\begin{equation}
\langle H\rangle=v\ne0.
\end{equation}
場をその真空値と揺らぎの和として書くと、
\begin{equation}
H=v+h.
\end{equation}
すると
\begin{equation}
H\phi^2=v\phi^2+h\phi^2.
\end{equation}
第一項は\(\phi\)について二次であり、その質量係数に寄与する。第二項は揺らぎの場\(h\)との相互作用として残る。\editorialnote{出典は、非零の真空値が質量係数になることを述べつつ文を終えている。示された代入は直接の代数的補完であり、その規格化と符号は完全なラグランジアンの規約に依存する。}

\begin{classroomqa}
\lecturetimestamp{02:01:29}
\audiencequestion{標準模型は仮想粒子の生成と消滅を含むのか。}
\lecturerresponse{もちろんである。仮想的な生成と消滅は、標準模型のラグランジアンに符号化された繰り返しの量子過程に本質的に備わっている。}
\end{classroomqa}
